\documentclass[10pt,a4paper]{amsart}
\usepackage{amsmath,amssymb,amsthm,graphicx}
\usepackage{verbatim}
\newenvironment{reference}{\comment}{\endcomment}
%\newenvironment{reference}{}{}
\newenvironment{slogan}{\comment}{\endcomment}
\newenvironment{history}{\comment}{\endcomment}

% For commutative diagrams we use Xy-pic
\usepackage[all]{xy}

% We use 2cell for 2-commutative diagrams.
\xyoption{2cell}
% Theorem environments.
%
\theoremstyle{plain}
\newtheorem{theorem}[subsection]{Theorem}
\newtheorem{proposition}[subsection]{Proposition}
\newtheorem{lemma}[subsection]{Lemma}

\theoremstyle{definition}
\newtheorem{definition}[subsection]{Definition}
\newtheorem{example}[subsection]{Example}
\newtheorem{exercise}[subsection]{Exercise}
\newtheorem{situation}[subsection]{Situation}

\theoremstyle{remark}
\newtheorem{remark}[subsection]{Remark}
\newtheorem{remarks}[subsection]{Remarks}

\numberwithin{equation}{subsection}

% Macros
%
\def\lim{\mathop{\mathrm{lim}}\nolimits}
\def\colim{\mathop{\mathrm{colim}}\nolimits}
\def\Spec{\mathop{\mathrm{Spec}}}
\def\Hom{\mathop{\mathrm{Hom}}\nolimits}
\def\Ext{\mathop{\mathrm{Ext}}\nolimits}
\def\SheafHom{\mathop{\mathcal{H}\!\mathit{om}}\nolimits}
\def\SheafExt{\mathop{\mathcal{E}\!\mathit{xt}}\nolimits}
\def\Sch{\mathit{Sch}}
\def\Mor{\mathop{\mathrm{Mor}}\nolimits}
\def\Ob{\mathop{\mathrm{Ob}}\nolimits}
\def\Sh{\mathop{\mathit{Sh}}\nolimits}
\def\NL{\mathop{N\!L}\nolimits}
\def\CH{\mathop{\mathrm{CH}}\nolimits}
\def\proetale{{pro\text{-}\acute{e}tale}}
\def\etale{{\acute{e}tale}}
\def\QCoh{\mathit{QCoh}}
\def\Ker{\mathop{\mathrm{Ker}}}
\def\Im{\mathop{\mathrm{Im}}}
\def\Coker{\mathop{\mathrm{Coker}}}
\def\Coim{\mathop{\mathrm{Coim}}}

% Boxtimes
%
\DeclareMathSymbol{\boxtimes}{\mathbin}{AMSa}{"02}

%
% Macros for moduli stacks/spaces
%
\def\QCohstack{\mathcal{QC}\!\mathit{oh}}
\def\Cohstack{\mathcal{C}\!\mathit{oh}}
\def\Spacesstack{\mathcal{S}\!\mathit{paces}}
\def\Quotfunctor{\mathrm{Quot}}
\def\Hilbfunctor{\mathrm{Hilb}}
\def\Curvesstack{\mathcal{C}\!\mathit{urves}}
\def\Polarizedstack{\mathcal{P}\!\mathit{olarized}}
\def\Complexesstack{\mathcal{C}\!\mathit{omplexes}}
% \Pic is the operator that assigns to X its picard group, usage \Pic(X)
% \Picardstack_{X/B} denotes the Picard stack of X over B
% \Picardfunctor_{X/B} denotes the Picard functor of X over B
\def\Pic{\mathop{\mathrm{Pic}}\nolimits}
\def\Picardstack{\mathcal{P}\!\mathit{ic}}
\def\Picardfunctor{\mathrm{Pic}}
\def\Deformationcategory{\mathcal{D}\!\mathit{ef}}
\usepackage{cleveref}
\crefname{originallemma}{original lemma}{original lemmas}
\Crefname{originallemma}{Original lemma}{Original lemmas}
\makeatletter
\newlabel{derived-lemma-essentially-constant-2-out-of-3}{{0G3A}{source}{}{sourcepointer-1}{}}
\newlabel{derived-lemma-essentially-constant-2-out-of-3@cref}{{[originallemma][0][]0G3A}{[1][0][]source}{}{}{}}
\makeatother
\crefname{originallemma}{original lemma}{original lemmas}
\Crefname{originallemma}{Original lemma}{Original lemmas}
\makeatletter
\newlabel{lemma-lift-map}{{0G2H}{source}{}{sourcepointer-2}{}}
\newlabel{lemma-lift-map@cref}{{[originallemma][0][]0G2H}{[1][0][]source}{}{}{}}
\makeatother
\crefname{originallemma}{original lemma}{original lemmas}
\Crefname{originallemma}{Original lemma}{Original lemmas}
\makeatletter
\newlabel{lemma-extension-by-zero}{{0G4L}{source}{}{sourcepointer-3}{}}
\newlabel{lemma-extension-by-zero@cref}{{[originallemma][0][]0G4L}{[1][0][]source}{}{}{}}
\makeatother
\crefname{originalreference}{original reference}{original references}
\Crefname{originalreference}{Original reference}{Original references}
\makeatletter
\newlabel{categories-remark-pro-category-copresheaves}{{0G2X}{source}{}{sourcepointer-4}{}}
\newlabel{categories-remark-pro-category-copresheaves@cref}{{[originalreference][0][]0G2X}{[1][0][]source}{}{}{}}
\makeatother
\title{The property of being pro-isomorphic to a Deligne system satisfies 2-out-of-3 along distinguished triangles}
\author{The Stacks Project Authors}
\date{stacks-project:duality.tex:local-snapshot-20261012; GNU FDL 1.2 or later}
\begin{document}
\maketitle
\noindent\textit{Editorial note.} All delivered mathematical text is the author's own, copied from the pinned Stacks Project chapter file duality.tex; no statement or proof step was rewritten, reordered or removed.\par\medskip
\noindent\textit{Editorial note.} Preamble provenance: the pinned chapter file opens with an input of the shared preamble and carries no preamble text of its own, so the project's single Stacks-wide preamble staged at sources/stacks-duality/preamble.tex is excerpted verbatim under the recipe fixed by LANE\_BRIEF 22.1, as three spans, lines 8-13 (verbatim package and the reference/slogan/history environments), lines 14-18 (Xy-pic setup) and lines 154-217 (the native theorem declarations and mathematical macros). Nothing was added, removed or redefined; the wrapper class and the amsmath/amssymb/amsthm/graphicx base packages are supplied by the assembler.\par\medskip
\noindent\textit{Editorial note.} The counted statement and proof contain no cross reference, no citation and no included file, so the delivered body is self contained.\par\medskip
\noindent\textit{Editorial note.} Canonical location of the counted proposition: Stacks tag 0G4Q (sources/stacks-duality/pinned-original-source.tex line 8236).\par\medskip
\noindent\textit{Allowed standard prerequisites (editorial).}\begin{itemize}
\item Essential constancy satisfies 2-out-of-3 along distinguished triangles. Derived Categories, Lemma labelled derived-lemma-essentially-constant-2-out-of-3, cited at proof lines 8245-8248. Gives essential constancy of the restriction of (L\_n) to U.
\item The colimit of Hom out of a Deligne system computes Hom on the open. Same chapter, Lemma labelled lemma-lift-map, cited at proof lines 8270-8272. Makes all but the middle vertical arrow of the diagram an isomorphism.
\item Existence of a Deligne system with prescribed constant value, and the characterisation of pro-isomorphism by colimits of Hom. Same chapter, Lemma labelled lemma-extension-by-zero, cited at proof lines 8278-8280, and Categories, Remark labelled categories-remark-pro-category-copresheaves, cited at proof lines 8280-8281. Concludes that (L\_n) is pro-isomorphic to a Deligne system.
\end{itemize}\medskip
\ifdefined\hypertarget\hypertarget{sourcepointer-1}{}\fi
\noindent\textit{Source reference.} Original lemma 0G3A: sources/stacks-chow/tags.txt, source lines 20085--20085. Cross-chapter standard prerequisite located in the project tag list staged in the local snapshot; invoked at proof line 8251 of the counted proof. It is retained only as a navigation pointer; every load-bearing step is inside the delivered span.\par\medskip
\ifdefined\hypertarget\hypertarget{sourcepointer-2}{}\fi
\noindent\textit{Source reference.} Original lemma 0G2H: sources/stacks-duality/pinned-original-source.tex, source lines 7920--7920. Same-chapter pointer into the author environment lemma at source line 7920; invoked at proof line 8270 of the counted proof. It is retained only as a navigation pointer; every load-bearing step is inside the delivered span.\par\medskip
\ifdefined\hypertarget\hypertarget{sourcepointer-3}{}\fi
\noindent\textit{Source reference.} Original lemma 0G4L: sources/stacks-duality/pinned-original-source.tex, source lines 8103--8103. Same-chapter pointer into the author environment lemma at source line 8103; invoked at proof line 8276 of the counted proof. It is retained only as a navigation pointer; every load-bearing step is inside the delivered span.\par\medskip
\ifdefined\hypertarget\hypertarget{sourcepointer-4}{}\fi
\noindent\textit{Source reference.} Original reference 0G2X: sources/stacks-chow/tags.txt, source lines 20072--20072. Cross-chapter standard prerequisite located in the project tag list staged in the local snapshot; invoked at proof line 8280 of the counted proof. It is retained only as a navigation pointer; every load-bearing step is inside the delivered span.\par\medskip
\begin{lemma}
\label{lemma-deligne-system-2-out-of-3}
Let $j : U \to X$ be an open immersion of Noetherian schemes.
Let
$$
K_n \to L_n \to M_n \to K_n[1]
$$
be an inverse system of distinguished triangles in
$D^b_{\textit{Coh}}(\mathcal{O}_X)$. If $(K_n)$ and $(M_n)$
are pro-isomorphic to Deligne systems, then so is $(L_n)$.
\end{lemma}

\begin{proof}
Observe that the systems $(K_n|_U)$ and $(M_n|_U)$ are essentially constant
as they are pro-isomorphic to constant systems.
Denote $K$ and $M$ their values. By Derived Categories, Lemma
\ref{derived-lemma-essentially-constant-2-out-of-3}
we see that the inverse system $L_n|_U$ is essentially constant as well.
Denote $L$ its value.
Let $N \in D^b_{\textit{Coh}}(\mathcal{O}_X)$. Consider the commutative
diagram
$$
\xymatrix{
\ldots \ar[r] &
\colim \Hom_X(M_n, N) \ar[r] \ar[d] &
\colim \Hom_X(L_n, N) \ar[r] \ar[d] &
\colim \Hom_X(K_n, N) \ar[r] \ar[d] &
\ldots \\
\ldots \ar[r] &
\Hom_U(M, N|_U) \ar[r] &
\Hom_U(L, N|_U) \ar[r] &
\Hom_U(K, N|_U) \ar[r] &
\ldots
}
$$
By Lemma \ref{lemma-lift-map} and the fact that isomorphic ind-systems
have the same colimit, we see that the vertical arrows two to the right
and two to the left of the middle one are isomorphisms. By the 5-lemma
we conclude that the
middle vertical arrow is an isomorphism. Now, if $(L'_n)$ is a Deligne system
whose restriction to $U$ has constant value $L$ (which
exists by Lemma \ref{lemma-extension-by-zero}), then
we have $\colim \Hom_X(L'_n, N) = \Hom_U(L, N|_U)$ as well.
Hence the pro-systems $(L_n)$ and $(L'_n)$ are
pro-isomorphic by
Categories, Remark \ref{categories-remark-pro-category-copresheaves}.
\end{proof}

\end{document}
